\chapter{讨论、局限性与可复现性}

\section{固定协议下结果的含义}

第~4~章结果表明，CLDNN 在本文固定协议中取得最高的整体准确率和宏平均 F1，并在已登记配对检验中优于 \path{resnet1d} 与 \path{iq_param_matched}。这一现象可以从结构上理解：CLDNN 先用卷积层提取局部时域模式，再通过循环单元建模序列依赖，相比单纯一维卷积基线具有更强的时序表征能力。与此同时，CLDNN 的优势仍然是协议内观察，不能扩展为跨数据集或跨信道条件的普遍结论。

固定划分的价值在于提高模型间比较的可比性。所有模型使用相同测试样本、相同训练种子集合和相同评价脚本，因而结果差异更容易归因于当前模型实现与训练配置，而不是测试集变化。其限制也同样明显：单一固定划分不能替代多划分重复实验，也不能证明模型对其他数据集或真实无线环境具有相同表现。

\section{STFT 融合的低信噪比权衡}

静态 I/Q--STFT 融合的结果说明，多视图输入的价值具有条件性。相对参数匹配 I/Q 控制模型，静态融合在低信噪比范围内有小幅正向差异；相对 CLDNN，该差异没有统计支持；在整体测试集上，静态融合又低于参数匹配 I/Q 控制模型。因此，本文只能得出“STFT 视图在低信噪比样本上可能提供局部补充信息”的保守结论。

这种权衡具有合理机制解释。低信噪比下，原始 I/Q 序列中的局部相位和幅度结构容易被噪声扰动，STFT 展开的局部频率能量分布可能提供另一种观察角度。但在中高信噪比样本上，I/Q 序列已经包含较清晰的时域判别线索，额外的时频分支和融合分类器可能引入训练难度、特征冗余或视图不一致，从而造成整体准确率下降。由此可见，融合结构的关键问题不是“是否增加视图”，而是如何在低信噪比收益与中高信噪比保持之间取得平衡。

\section{门控融合负结果的意义}

门控融合没有改善静态融合，是本文值得保留的负结果。该模型的设计动机是用样本级标量门控自动调节 I/Q 与 STFT 表征贡献，但当前结果显示，门控权重并未带来更好的整体或低信噪比表现。可能原因包括：门控输入没有显式信噪比信息，标量权重难以表达复杂的类别和信噪比依赖关系；门控子网络也可能在有限训练协议下增加优化难度。

这一负结果不应被解释为所有注意力或门控机制都无效。它只说明本文登记的 \path{gated_fusion_iq_stft}、当前实现和当前训练协议没有形成预期收益。后续若研究显式信噪比感知门控、类别条件门控或更稳定的多视图训练策略，应以新的协议独立评估，不能把设想直接写成本报告结论。

\section{MCLDNN 异常的解释边界}

MCLDNN 的两个训练种子出现近似随机水平的单类预测，说明该结构在当前实现和训练配置下存在明显稳定性问题。保留异常种子并不是降低报告质量，而是维护固定协议的一致性：如果只删除失败种子，主表会变得更好看，但不再忠实反映登记实验矩阵。

该异常更适合解释为优化或初始化敏感性，而不是直接证明数据错误或模型家族无效。已有审计材料没有提供删除种子的证据，因此本文保留完整三种子聚合。后续若要修复或诊断 MCLDNN，应使用新的模型标识、独立输出目录和完整统计检验，不能用诊断性结果替换本文主表。

\section{复杂度与运行时证据限制}

本文的复杂度讨论只覆盖参数量、训练时间背景和受控 CUDA 前向延迟。该测量能帮助读者理解不同模型在相同硬件和相同计时设置下的前向耗时差异，但不能代表完整部署成本。尤其对于 STFT 和 I/Q--STFT 模型，当前延迟表不包含完整时频预处理、数据搬运、CPU 推理或端到端流水线时间。

因此，本文不作“更适合部署”“更高计算效率”或“硬件可迁移性能更好”的判断。若未来需要系统级效率分析，应同时测量 CPU/GPU、不同批量大小、预处理耗时、内存峰值、吞吐率和端到端延迟，并将该协议与本报告的受控前向延迟区分开。

\section{可复现性与证据边界}

本文可复现性主要来自三个层面：固定划分和三训练种子保证比较口径一致；测试集预测归档支持同一样本上的配对统计检验；聚合表、低信噪比分组表和受控延迟表使正文数值可追溯。表~\ref{tab:course-artifact-boundary} 用课程论文口径概括这些材料的用途，具体路径、环境和同步状态放在附录 A。

\begin{table}[htbp]
  \centering
  \caption{课程报告使用的主要证据材料及边界}
  \label{tab:course-artifact-boundary}
  \begin{tabularx}{\textwidth}{lXX}
    \toprule
    材料类型 & 主要作用 & 使用边界 \\
    \midrule
    聚合结果表 & 支撑主结果、低/中/高信噪比分组和复杂度背景。 & 只用于 \datasetname{} 固定协议，不外推到其他数据集。 \\
    测试集预测归档 & 支撑配对自助法和 McNemar 检验。 & 只覆盖 9 个模型行、3 个训练种子和已登记模型对。 \\
    受控延迟表 & 支撑参数量与 CUDA 前向延迟讨论。 & 不代表 CPU、预处理完整耗时或端到端部署成本。 \\
    环境与同步记录 & 说明本地写作环境、服务器实验环境和权重归档缺口。 & 不等同于本地完整训练权重归档。 \\
    \bottomrule
  \end{tabularx}
\end{table}

附录图~\ref{fig:evidence-boundary} 进一步说明证据标签与排除规则。诊断性或工程检查材料可以用于说明后续工作边界，但不能进入主结果、统计检验、摘要和结论。当前本地材料足以支撑报告写作和统计解释，但并不包含完整 27 个训练权重文件；若未来需要模型加载或权重级复现，应先进行独立权重同步和哈希核验。
